\documentclass[11pt, letterpaper]{article}
\usepackage[utf8]{inputenc}
\usepackage{geometry}
\usepackage{fancyhdr}
\usepackage{amsmath, amssymb}
\usepackage{xcolor}
\usepackage{listings}
\usepackage{enumitem}
\usepackage{hyperref}

% Make parindent 0
\setlength{\parindent}{0pt}
\setlist{itemsep=0.5em}
\setlength{\parskip}{0.5em}

\geometry{margin=1in}
\setlength{\headheight}{14pt}
\pagestyle{fancy}
\fancyhf{}
\lhead{\textbf{Student Name:} \textcolor{red}{YOUR NAME}} % CHANGE NAME HERE
\rhead{\textbf{Data Analysis and Regression, Homework 1}}
\cfoot{\thepage}

\definecolor{codegreen}{rgb}{0,0.6,0}
\definecolor{codegray}{rgb}{0.5,0.5,0.5}
\definecolor{codepurple}{rgb}{0.58,0,0.82}
\definecolor{backcolour}{rgb}{0.95,0.95,0.92}

\lstdefinestyle{mystyle}{
    backgroundcolor=\color{backcolour},   
    commentstyle=\color{codegreen},
    keywordstyle=\color{magenta},
    numberstyle=\tiny\color{codegray},
    stringstyle=\color{codepurple},
    basicstyle=\ttfamily\footnotesize,
    breakatwhitespace=false,         
    breaklines=true,                 
    captionpos=b,                    
    keepspaces=true,                 
    numbers=left,                    
    numbersep=5pt,                  
    showspaces=false,                
    showstringspaces=false,
    showtabs=false,                  
    tabsize=2
}
\lstset{style=mystyle}

\newcommand{\problem}[1]{\section*{Problem #1}}
\newcommand{\subproblem}[1]{\subsection*{Problem #1}}
\newcommand{\answer}{\textbf{\textit{\textcolor{red}{Answer:}}}\\}

\begin{document}

\section*{Instructions}
\begin{itemize}
    \item Homework 1 is due Monday, July 27th at 18:00 Chicago Time.
    \begin{itemize}
        \item We will not accept any submissions past 18:00:00, even if they are only one second late.
    \end{itemize}
    \item You \textbf{must} upload the following files to the class Canvas:
    \begin{itemize}
        \item \texttt{LASTNAME\_FIRSTNAME.pdf}
        \item \texttt{LASTNAME\_FIRSTNAME.ipynb}
    \end{itemize}
    \item Your code notebook \textbf{must} be runnable using my environment outlined in class 1 (Python 3.14, and the \texttt{requirements.txt}).
    \item You \textbf{must} use this template file and fill out your solutions for the written portion.
    \item Please note that your last name and first name should match what you appear on Canvas as.
    \item Include code snippets where required, as well as math and equations.
    \item Be \textit{concise} where possible, all of the homework problems can be answered in a few lines of math, code, and words.
\end{itemize}
\hrule
\vspace{0.5cm}

\pagebreak

\problem{1: One-Dimensional Data}

Load in the data from the GitHub repository for this class.
\begin{lstlisting}[language=python]
import pandas as pd

df = pd.read_csv(
    "https://raw.githubusercontent.com/tobiasdelpozo/data-analysis-2026/refs/heads/master/homework/homework_1/questions/homework_1_data.csv"
)
\end{lstlisting}

\subproblem{1.1}

For the feature labeled \texttt{X1}, compute the mean, median, variance, and standard deviation. 
Report your numbers below (rounded to at least 4 decimal places).

\answer

\subproblem{1.2}

Display a histogram of the feature \texttt{X1} using 50 bins. Do you think that 
the statistics you computed in 1.1 are good descriptors of the data? Include the graph below,
and explain your reasoning in 1-2 sentences. \footnote{Hint: you can use \texttt{\textbackslash includegraphics\{\}} to include images in \LaTeX.}

\answer

\subproblem{1.3}

Using the same feature \texttt{X1}, come up with some metrics that 
are descriptive of the distribution of the data. Note, this is open-ended,
so think about what the data looks like, and how a human would describe it.

\answer


\pagebreak

\problem{2: kNN Regression}

This problem uses the same dataset as Problem 1.

We're going to implement a k-Nearest Neighbors regression model. Unless otherwise specified, use an 80/20 train/test split for all parts of this problem.

\subproblem{2.1}

Display a plot of \texttt{X2} versus the \texttt{target} variable. What 
do you notice about the relationship between these two variables?

\answer

\subproblem{2.2}

Implement a kNN regression model from scratch. You may use \texttt{numpy} and \texttt{pandas}, but you may not use any machine learning libraries (e.g. \texttt{scikit-learn}).

Your model should take in 4 parameters:
\begin{itemize}
    \item \texttt{X\_train}: training features
    \item \texttt{y\_train}: training target variable
    \item \texttt{X\_test}: testing features
    \item \texttt{k}: number of neighbors to use
\end{itemize}
And it should output the predicted values for \texttt{X\_test}.

The algorithm you should use is as follows:
\begin{enumerate}
    \item For each test point, compute the Euclidean distance to all training points.
    \item Identify the k-nearest neighbors based on these distances.
    \item Compute the predicted value as the mean of the target variable of these k-nearest neighbors.
\end{enumerate}
Note that we are only considering a single feature for this problem, so the Euclidean distance is
simply $\sqrt{(x_{\text{test}} - x_{\text{train}})^2}$.

Include your code implementation below.\footnote{
    Hint: you can use \texttt{\textbackslash lstlisting[language=python]} to include Python code snippets.
}

\answer

\subproblem{2.3}

Randomly split the data into training and testing sets (80/20 split), and
report the Mean Squared Error (MSE) of your kNN regression model on the test set for $k=5$.

\answer

\subproblem{2.4}

For $k \in \{1, 5, 10, 20, 50, 100 \}$, compute the MSE on the test set and plot the results (k values on the x-axis, MSE on the y-axis).
What value of $k$ gives the best performance on the test set?

\answer


\subproblem{2.5}

Which value of $k$ do you think has the highest bias? And which has the highest variance? Explain your reasoning in 1-2 sentences.

\answer

\problem{3: Linear Regression}

Using the same dataset as Problems 1 and 2, we are going to explore linear regression.

\subproblem{3.1}

Using \texttt{statsmodels}, fit a linear regression model to predict \texttt{y} using \texttt{X3}.
You should \textbf{not} use an intercept term in your model. Report your
$\beta$ coefficient below:

\answer

\subproblem{3.2}

Re-run the linear regression model from 3.1, but this time include an intercept term.
What are your new $\beta$ coefficients (intercept and slope)?

\answer

\subproblem{3.3}

Do the following data transformations:
$$\tilde{y} = y - \bar{y} \quad \quad \tilde{X}_3 = X_3 - \bar{X}_3$$

Re-run the linear regression model using $\tilde{y}$ and $\tilde{X}_3$, without an intercept term.
What is your $\beta$ coefficient? How does it compare to your answer in 3.1?

\answer

\subproblem{3.4}

Inspect your data. Display a scatter plot of \texttt{X3} versus \texttt{y}. What do you notice about the relationship between these two variables?
Is a linear model appropriate for this data? Explain your reasoning in 1-2 sentences.

\answer

\subproblem{3.5}

Define a new feature $\texttt{X3\_sin}$ as follows: 
$$\texttt{X3\_sin} = \sin(\texttt{X3})$$

Fit a linear regression model to predict \texttt{y} using \texttt{X3\_sin}, including an intercept term.
Report your $\beta$ coefficients (intercept and slope) below:

\answer

\subproblem{3.6}

Display a plot of \texttt{X3\_sin} versus \texttt{y}. Do you think a linear model is appropriate for this data? Explain your reasoning in 1-2 sentences.

\answer

\end{document}